% Source: https://artofproblemsolving.com/community/c3h2653643_math_puzzle_21082021
% Stable identity: aops-topic-2653643-q01
The six equations are
\[\dfrac{AB}{C}=24,\quad D+EF=48,\quad G-H+I=4,\quad \dfrac AD-G=-5,\quad BE-H=36,\quad C+FI=11.\]
\begin{itemize}
\item From $\dfrac AD=G-5$: $\dfrac AD$ is a positive integer, not $1$ (as $A\ne D$), and at most $9-5=4$. So $\dfrac AD\in\{2,3,4\}$ and $G\ge7$.
\item From $D+EF=48$ with $D\le9$: $39\le EF\le47$, and the products of two different digits in this range are $40=5\cdot8$, $42=6\cdot7$, $45=5\cdot9$, giving $D=8,6,3$. But $A\ge2D$ rules out $D=8$ and $D=6$, so $D=3$ and $\{E,F\}=\{5,9\}$.
\item Then $A=3\cdot\dfrac AD\in\{6,9,12\}$; $12$ is not a digit and $9$ is taken by $E$ or $F$, so $A=6$ and $G=\dfrac63+5=7$.
\item $AB=24C$ gives $B=4C$, so $(C,B)=(1,4)$ or $(2,8)$.
\item $C+FI=11$: if $F=9$ then $I=1$, $C=2$, $B=8$, $E=5$; if $F=5$ then $(I,C)=(1,6)$ or $(2,1)$, but $C=6=A$ is impossible, so $C=1$, $B=4$, $E=9$, $I=2$.
\item $H=BE-36$: in the second case $H=36-36=0$, impossible. In the first, $H=40-36=4$.
\end{itemize}
So $\boxed{A=6,\ B=8,\ C=2,\ D=3,\ E=5,\ F=9,\ G=7,\ H=4,\ I=1}$, which uses each digit once. Check: $6\times8\div2=24$, $3+5\times9=48$, $7-4+1=4$, $6\div3-7=-5$, $8\times5-4=36$, $2+9\times1=11$.
